\subsection*{MT-M22. The fixed Hecke sector of Meyer's adelic zero-space}

For every fixed primitive unitary Hecke \(L(s,\eta)\), the
canonical quotient \(Q_\eta\) of MT-M18 is isomorphic to the
corresponding fixed compact-character sector of Meyer's
idele-class zero-space.  The isomorphism is unconditional
and compatible with convolution, including all zero
multiplicities.  The sector is first given its quotient
Fr\'echet topology from \(H_-^\chi\); the proof also
identifies the quotient bornology by an explicit splitting.

We fix the character convention explicitly.  Write
\[
 L(s,\eta)=\prod_{\mathfrak p}
       (1-c_{\mathfrak p}(N\mathfrak p)^{-s})^{-1},
\]
where \(c_{\mathfrak p}=0\) at a ramified place.  Let
\(\beta\) be the associated unitary idele-class character
in Meyer's convention: for his \(S\)-prime element
\(p_v\in C_S\), at every unramified \(v\notin S\),
\[
                      \beta(p_v)^{-1}=c_{\mathfrak p_v}.
                                                               \tag{M22.1}
\]
Thus the component is parameterized by
\(\omega_s=\beta|\cdot|^s\).
This convention is dictated by the inverse in Meyer's
Euler factor, rather than by a choice of the name
\(\eta\) or \(\eta^{-1}\) for its idele realization.
For example, \(L(s,\eta)=\zeta_K(s+i\tau)\) corresponds
to \(\beta=|\cdot|^{i\tau}\).
If an idele character has instead been labeled by
its direct values \(c_{\mathfrak p_v}\) on Meyer's
\(p_v\), it must be inverted to obtain \(\beta\).

Let \(C=C_K\), \(C^1=\ker|\cdot|\), and choose a
continuous smooth splitting \(a:\mathbb R\to C\)
with \(|a(t)|=e^t\).
Normalize Haar measure on \(C^1\) to one, so
\(d^\times x=dc\,dt\) in \(x=ca(t)\).
Put \(\chi=\beta|_{C^1}\).
Use Meyer's spaces and embeddings
\[
 H_-=\mathcal S(C)_{\mathbb R},\qquad
 H_+=\mathcal S(\mathbb A_K)/K^\times,\qquad
 H_-^0=H_-/(H_+\cap H_-),
\]
where the intersection uses his embeddings into
\(\mathcal S(C)_{]1,\infty[}\oplus
  \mathcal S(C)_{]-\infty,0[}\).
Superscript \(\chi\) denotes the subspace on which
left translation by \(c\in C^1\) acts by \(\chi(c)\).
Then
\[
 (\mathcal L_\beta f)(x)=\beta(x)^{-1}f(\log|x|)
\]
induces an isomorphism
\[
              Q_\eta \ \simeq\ (H_-^0)^\chi.                  \tag{M22.2}
\]
It is topological for the specified quotient topology.
Norm translation by \(t\) corresponds to idele
translation by \(a(t)\), with the explicit unitary
factor \(\beta(a(t))\).

The primary source is Ralf Meyer, \emph{On a representation
of the idele class group related to primes and zeros of
\(L\)-functions}, Duke Math. J.~127 (2005), 519--595,
author version \texttt{math/0311468v3}.
The exact input is the weighted spaces in Definition 4.1
and (22), author pp.~24--25; the embeddings in Lemma 5.4,
pp.~40--41; and the image criterion in \S5.7, pp.~46--47,
following Lemma 5.10.  Theorem 5.1, pp.~38--39, supplies
the finite-level description used there.
Theorems 5.8, pp.~43--44, and 5.11, p.~47, are the
classical summability and spectral-multiplicity statements.
The comparison below uses the image criterion,
rather than inferring an isomorphism from equal spectra.

\begin{proof}
\emph{The exact test-function sector.}
A norm splitting exists by varying a positive component
at one archimedean place with the appropriate local
norm exponent.  In the resulting product
\(C=C^1\times a(\mathbb R)\), every \(\chi\)-homogeneous
function has the form \(\chi(c)^{-1}u(t)\).
Since \(\beta\) is unitary,
\(\beta(a(t))=e^{i\tau t}\) for a fixed real \(\tau\).
Multiplication by this fixed modulation preserves
every exponential-Schwartz seminorm, in both
directions, by the Leibniz formula.

Consequently \(\mathcal L_\beta\) is a topological
isomorphism from \(\mathcal E\) onto \(H_-^\chi\).
Indeed the norm-direction invariant derivative becomes
\(d/dt\); derivatives in the compact direction act
on its one fixed smooth character by constants.
All norm weights are \(e^{\alpha t}\), and the
requirement for every real \(\alpha\) is equivalent
to the seminorms \(e^{N|t|}|f^{(j)}(t)|\).
Polynomial factors in \(t\) are absorbed by increasing
\(N\).  The character is trivial on a suitable
compact subgroup with elementary compact quotient,
so these functions occupy one fixed Bruhat level.
There is no limit over varying compact characters.

Haar normalization as in Meyer's \S3.1, p.~16,
gives exactly
\[
 \widehat{\mathcal L_\beta f}(\omega_s)
       =\int_C\beta(x)^{-1}f(\log|x|)
                      \beta(x)|x|^s\,d^\times x
       =\mathcal Mf(s).                                        \tag{M22.3}
\]
The transform is zero on all other compact-character
components by orthogonality on \(C^1\).
A direct convolution calculation, using
\(\beta(y)^{-1}\beta(y^{-1}x)^{-1}=\beta(x)^{-1}\),
shows
\[
  (\mathcal L_\beta f)*(\mathcal L_\beta g)
                         =\mathcal L_\beta(f*g).
\]
Also
\[
 \lambda_b\mathcal L_\beta f
   =\beta(b)\mathcal L_\beta
         \bigl(f(\,\cdot-\log|b|)\bigr),\qquad b\in C.            \tag{M22.4}
\]
These identities specify the representation convention.

\emph{Meyer's image condition on this component.}
Choose once and for all a sufficiently large finite
set \(S\), containing the archimedean places, all
ramification of \(\beta\), and the places needed for
his normalized additive character and differential
idele.  His multiplier \(\mathcal L_S\) has transform
\[
 \widehat{\mathcal L_S}(\omega)
       =\prod_{v\notin S}(1-\omega(p_v)^{-1})^{-1}.
\]
By (M22.1), on \(\omega_s\) this is the partial finite
Hecke function
\[
 L_\eta^S(s)
       =L(s,\eta)D_{\eta,S}(s),\qquad
 D_{\eta,S}(s)=
        \prod_{v\in S_{\mathrm{fin}}}
             (1-c_{\mathfrak p_v}(N\mathfrak p_v)^{-s}).         \tag{M22.5}
\]
A ramified factor in \(D_{\eta,S}\) is one.
The reflected quasi-character is
\[
 \omega_s^*=|\cdot|\omega_s^{-1}
                    =\beta^{-1}|\cdot|^{1-s},
\]
and therefore
\[
              \widehat{\mathcal L_S}(\omega_s^*)
                       =L_{\bar\eta}^S(1-s).                   \tag{M22.6}
\]
Here \(\bar\eta\) denotes the dual unitary Hecke
character, with conjugate Euler coefficients.

For an entire strip-Schwartz function \(F(s)\)
on this component, the equivalent reduced
image criterion printed on Meyer p.~47 is
\[
 \begin{array}{ll}
 F(\alpha+it)/L_\eta^S(\alpha+it)
        \text{ is Schwartz in }t,&0<\alpha\le1,\\[2mm]
 F(it)/L_{\bar\eta}^S(1-it)
        \text{ is Schwartz in }t.&
 \end{array}                                                    \tag{M22.7}
\]
All apparent zero singularities must be removable.
The source's remaining conditions for
\(\alpha>1\) and \(\alpha<0\) follow from \(F\)
being the transform of \(H_-\).
Zero extension to other compact-character
components causes no additional Schwartz condition.
Thus (M22.7) characterizes membership of
\(\mathcal L_\beta f\) in \(\Sigma(H_+)\),
where \(F=\mathcal Mf\).

This is also the relevant intersection in
the quotient, rather than only its first-coordinate
image.  Meyer's Poisson formula gives
\[
 \Sigma u-J\Sigma\mathfrak F^*u
                         =c_0-c_1|\cdot|^{-1}.
\]
The second term lies in
\(\mathcal S(C)_{]-\infty,0[}\).
If \(\Sigma u\in H_-\), so does the right side.
Its decay as \(|x|\to0\) forces \(c_0=c_1=0\):
these distinct powers cannot decay after every
negative weight.  Hence
\(i_+u=i_-(\Sigma u)\).
Conversely, an element of the intersection has
its first coordinate in \(\Sigma(H_+)\cap H_-\).
Compact-character projection commutes with the
maps.  Therefore (M22.7) characterizes
\((H_+\cap H_-)^\chi\).

\emph{Removing the image criterion's growth gate.}
If (M22.7) holds, every zero of \(\Xi_\eta\) in
\(0<\operatorname{Re}s\le1\) is a zero of \(F\)
with at least its full multiplicity, by its first
condition.  The omitted finite Euler factors have
no zeros in that right half-plane, and the gamma
factors have neither zeros nor poles there.
The possible pure-norm pole is cancelled in
\(\Xi_\eta\) and is not a zero of it.
If a nontrivial zero is on \(\operatorname{Re}s=0\),
the functional equation identifies its multiplicity
with that of \(L_{\bar\eta}(1-s)\); the second
condition gives the same conclusion.
Thus membership implies all the zero-jet
vanishings of MT-M18.

Conversely suppose \(F=\mathcal Mf\) vanishes
on the full divisor of \(\Xi_\eta\).
Set \(G=F/\Xi_\eta\), which is entire.
The rectangle proof of (M18.14) gives,
for every \(A,N>0\), the sharpened estimate
\[
 \sup_{|\sigma|\le A}|G(\sigma+it)|
       \le C_{A,N,f}e^{\kappa|t|}(2+|t|)^{-N},
 \qquad \kappa=\pi[K:\mathbb Q]/4.                              \tag{M22.8}
\]
Indeed the boundary lower bound for \(\Xi_\eta\)
costs \(e^{\kappa|n|}(2+|n|)^{B_A}\).
Use the decay power \(N+B_A\), rounded up and
increased if needed, for \(F\) on that boundary.
The maximum principle and bounded block heights
give (M22.8), with no zero-spacing hypothesis.

Write the standard completion as
\[
 \Xi_\eta(s)=P_\eta(s)q_\eta^{s/2}
                  \Gamma_\eta(s)L(s,\eta).
\]
On each closed strip in \(\operatorname{Re}s>0\),
the unitary archimedean gamma product is
holomorphic and bounded above by
\(e^{-\kappa|t|}\) times a fixed power
of \(2+|t|\), by Stirling's formula.
Consequently
\[
        F(s)/L(s,\eta)
          =G(s)P_\eta(s)q_\eta^{s/2}\Gamma_\eta(s)
\]
is holomorphic and rapidly decreasing there,
by (M22.8).  The reciprocal of a finite-\(L\)
pole is a zero, so pure-norm poles cause
no exception.

The completed functional equation gives
\[
 \Xi_\eta(s)
   =\epsilon_\eta P_\eta(s)q_\eta^{(1-s)/2}
            \Gamma_{\bar\eta}(1-s)L(1-s,\bar\eta),
                    \qquad|\epsilon_\eta|=1.
\]
It follows that \(F(s)/L(1-s,\bar\eta)\)
is holomorphic and rapidly decreasing
on each closed strip in
\(\operatorname{Re}s<1\).
The reflected gamma arguments have positive
real part there.  Thus neither half-plane
division asks \(F\) to cancel trivial
finite-\(L\) zeros outside the appropriate half.

For \(\delta>0\), every nontrivial omitted
Euler factor satisfies
\[
 |1-c_{\mathfrak p}(N\mathfrak p)^{-s}|
                         \ge1-(N\mathfrak p)^{-\delta}
       \quad(\operatorname{Re}s\ge\delta).
\]
The ramified factors are one.
Thus \(D_{\eta,S}^{-1}\) is bounded on each
such right half-plane, as is
\(D_{\bar\eta,S}(1-s)^{-1}\) when
\(\operatorname{Re}s\le1-\delta\).
Their derivatives are bounded on smaller
fixed strips.  Therefore
\[
 \begin{split}
 F/L_\eta^S&\text{ is strip-Schwartz on }\operatorname{Re}s>0,\\
 F/L_{\bar\eta}^S(1-s)&\text{ is strip-Schwartz on }
                                      \operatorname{Re}s<1.
 \end{split}                                                    \tag{M22.9}
\]
Cauchy's formula on slightly larger strips
gives every vertical derivative estimate.
These statements imply (M22.7).
We have proved
\[
 \mathcal L_\beta^{-1}\bigl((H_+\cap H_-)^\chi\bigr)
   =\{f:(\mathcal Mf)^{(j)}(\rho)=0,\
          \Xi_\eta(\rho)=0,\ j<m_\rho\}
   =\overline{\psi_\eta*\mathcal E},                             \tag{M22.10}
\]
where the last equality is MT-M18.

\emph{The precise quotient topology and bounded lifting.}
Give \(H_-^\chi\simeq\mathcal E\) its
just-identified Fr\'echet topology and its
bounded-set bornology.  At this fixed character
these are precisely the source's weighted
Schwartz structure: it lies at one Bruhat
level, and the countable family of real
weight endpoints in (22) generates its
seminorms and bounded sets.
Here bounded and precompact sets coincide,
so the convention of using the precompact
bornology gives the same result.  Indeed,
for a bounded sequence in \(\mathcal E\),
the bounds on one additional derivative
give equicontinuity of every derivative
on each compact interval.  Diagonal
Arzel\`a--Ascoli gives a subsequence
converging with every derivative on every
compact interval.  For a prescribed
\(p_{N,j}\), the uniform \(p_{N+1,j}\)
bound makes the weighted tails outside
\([-R,R]\) at most \(Ce^{-R}\).
The subsequence is therefore Cauchy in
every defining seminorm and converges
in the complete space \(\mathcal E\).
This proves relative compactness of every
bounded set in its metrizable topology;
the converse is immediate.
The subspace in (M22.10) is closed.
Hence (M22.10) immediately gives a
topological isomorphism from \(Q_\eta\)
onto the sector quotient furnished with
the quotient topology from \(H_-^\chi\).

Here is why the original quotient bornology
agrees as well.  The maps \(i_+\) and \(i_-\)
are bornological embeddings by the cited
source.  Therefore the bornology of their
intersection equals the subspace bornology
inherited from \(H_-^\chi\): a subset of the
intersection is bounded in either embedded
space exactly when it is bounded in the
common ambient space.  No topology on
\(H_+^\chi\), or on its quotient, is being
assumed in this argument.

Compact-character averaging is a bounded
projection preserving \(H_+\) and \(H_-\).
At each finite level it is ordinary integration
over a compact group against its smooth
character, and these projections agree under
the level inclusions.
It preserves their intersection and splits
both spaces into its image and kernel.
The algebraic quotient and its quotient
bornology therefore split accordingly.
In particular \((H_-^0)^\chi\) is the
bornological quotient
\(H_-^\chi/(H_+\cap H_-)^\chi\);
representatives can be projected before
passing to the quotient.

Finally the continuous linear section
\(S:s(J)\to\mathcal E\) in MT-M18 supplies
an explicit lifting map for this quotient,
after transport by \(\mathcal L_\beta\).
Every bounded set in the quotient
Fr\'echet space lifts to a bounded set
through this section.  Conversely the
quotient map is continuous, so it takes
bounded sets to bounded sets.
The same two assertions hold for
precompact sets, using continuity of
the section and quotient map.
The quotient itself has bounded sets
precompact, since its continuous section
maps them to bounded, hence precompact,
subsets of \(\mathcal E\), and its
quotient map sends these back to the
original sets.
Thus the source's quotient bornology
is exactly the bornology of the specified
Fr\'echet quotient.  This is an explicit
splitting argument, not an inference
from equality of spectra.
Together with (M22.3)--(M22.4) and
multiplicativity, it proves all parts
of (M22.2).
\end{proof}

The stable blocks and nuclear traces in
MT-M18--MT-M19 therefore apply to this
independently defined adelic summation
quotient in each fixed Hecke sector.
All constants may depend on the character
and its infinity type; no uniform
decomposition over all characters is
asserted.  The trace positivity and
faithful-operator distinctions of
MT-M19--MT-M20 retain their exact GRH
and multiplicity hypotheses.
This identifies Meyer's specified
quotient, rather than an unspecified
quotient in Masullo's draft.
Its geometric location is the norm
direction of the idele-class action,
Tate's Poisson summation, and cote~1's
nuclear-trace framework as carried in
SA23 geo-H-lidskii/geo-H-genus.
The primary checks cover the named
normalization, image, and embedding
passages; no complete-paper audit
is claimed.

